% ==============================================================================
% CAPÍTULO 6 - CONSIDERAÇÕES FINAIS
% ==============================================================================

\chapter{Considerações Finais}
\label{cap:conclusao}

A realização do estágio na FK Oficina permitiu uma análise profunda sobre as complexidades inerentes ao desenvolvimento de plataformas de larga escala no modelo SaaS. O projeto Oficina Master consolidou-se como uma peça fundamental na arquitetura da empresa, transformando processos que antes eram puramente operacionais em ativos tecnológicos automatizados.

Do ponto de vista técnico, a implementação do sistema multi-tenancy com isolamento de banco de dados validou a importância da aplicação rigorosa de padrões de projeto (\textit{Design Patterns}) como \textit{Singleton}, \textit{Dependency Injection} e \textit{Observer}. A transição para o Laravel 12 e Vue.js 3 demonstrou como a escolha de tecnologias modernas e bem suportadas pela comunidade pode acelerar o desenvolvimento sem comprometer a manutenibilidade do código.

Os resultados obtidos transcendem a mera entrega de software; houve uma otimização mensurável na eficiência operacional da equipe de suporte técnico e uma redução drástica nos erros de sincronização tributária. Além disso, a integração com o Mercado Pago pavimentou o caminho para a escalabilidade financeira do ecossistema, automatizando o ciclo de vida das licenças desde a prospecção até a renovação.

Como conclusão, este estágio foi determinante para o amadurecimento profissional, unindo a fundamentação acadêmica da Ciência da Computação com as exigências de alta performance e disponibilidade do mercado de TI. Para evoluções futuras, vislumbra-se a implementação de inteligência artificial para monitoramento preditivo de anomalias nos bancos de dados dos \textit{tenants} e a migração de partes críticas do sistema para uma arquitetura \textit{serverless}, buscando otimizar ainda mais os custos de infraestrutura.
